\scanpage{85}

quand la suite $(a_{\sigma i})$ est décroissante, on est donc \deleted{\unsure{en réarrangeant au besoin les indices|on peut donc se ramener au besoin}} \deleted{\unsure{supposer, pour prouver (8), que la suite} $a_1\geq a_2\ldots\geq a_n$} ramené à prouver que si $a_1\geq a_2\ldots\geq a_n$, on a $\sum a_i\alpha_i\leq\deleted{\unsure{a_i\alpha|\sum a_i\alpha_i}}\sum a_i\beta_i\deleted{\unsure{\sum a_i\alpha|\sum a_i\alpha_i}}$, \deleted{\unsure{en supposant ou que}}\inserted{pourvu que} les $a_i$ soient positifs \deleted{si on n'est pas dans le cas b) du lemme} ou qu'on ait $\sum\alpha_i=\sum\beta_i$. Or on a en effet

\[
\begin{aligned}
\sum a_i\alpha_i={}&(a_1-a_2)\alpha_1+(a_2-a_3)(\alpha_1+\alpha_2)+\cdots+(a_{n-1}-a_n)(\alpha_1+\cdots+\alpha_{n-1})\\
&+a_n(\alpha_1+\cdots+\alpha_n)
\end{aligned}\tag*{(9)}
\]

\deleted{\unsure{et on a}} et une formule analogue pour $\deleted{\unsure{\sum a_i\alpha_i}}\sum a_i\beta_i$. Comme les coefficients $(a_1-a_2),\ldots,(a_i-a_{i\unsure{-1|+1}}),\ldots(a_{n-1}-a_n)$ sont positifs, \deleted{\unsure{on conclut que la somme relative|ainsi que}} et que $\deleted{\unsure{\sum}}\alpha_1\leq\beta_1,\ldots(\alpha_1+\cdots+\alpha_{n-1})\leq(\beta_1+\cdots+\beta_{n-1})$, on en conclut que tous les termes de la somme relative aux $\alpha_i$ sont majorés par les termes correspondants de la somme relative aux $\beta_i$, à l'exception \inserted{à priori} du dernier. Mais on a aussi $a_n(\alpha_1\unsure{+\cdots+|\ldots}\alpha_n)\leq a_n(\beta_1+\cdots+\beta_n)$ dans \deleted{\unsure{du terme du dernier}}\inserted{le cas a)} du lemme ($a_n\geq0$, $\sum\alpha_i\leq\sum\beta_i$) et dans le cas b) ($\sum\alpha_i=\sum\beta_i$), d'où \deleted{la conclusion}\inserted{l'inégalité (1)}.

\deleted{\unsure{Notons}} Signalons en passant le

% slanted
\margin{\uline{petits caractères}}

\uline{Corollaire \unsure{bis|1bis|1a}} Sous les conditions de \deleted{la prop. \unsure{du} corollaire 1, b,} \inserted{th. 1, 2°, b)} \deleted{\unsure{(hypothèse du corollaire 1, et|(resp. du corollaire 1, a))} \unsure{la prop. corollaire, b)|les mêmes hypothèses}} on aura aussi, \deleted{\unsure{si on a pour tout} \unsure{$\sigma$}}

\[
f\left(\frac1{\sigma_1},\ldots\frac1{\sigma_n}\right)\leq f\left(\frac1{\rho_1},\ldots,\frac1{\rho_n}\right)\tag*{(10)}
\]

En effet, de $\rho_1\ldots\rho_n=\sigma_1\ldots\sigma_n$ on tire pour tout \inserted{$1\leq\unsure{M|N}\leq n$ :}

\[
\rho_M\rho_{M+1}\ldots\rho_n=\frac{\sigma_1\ldots\sigma_{M-1}}{\rho_1\ldots\rho_{M-1}}\sigma_M\ldots\sigma_n,\qquad\text{d'où }\rho_M\ldots\rho_n\leq\sigma_M\ldots\sigma_n
\]

en vertu de $\sigma_1\ldots\sigma_{M-1}\leq\rho_1\ldots\rho_{M-1}$. Il s'ensuit que que les suites $(\frac1{\rho_n},\ldots\frac1{\rho_1})$ et $(\unsure{(}\frac1{\sigma_n},\ldots\frac1{\sigma_1})$ satisfont encore aux conditions du corollaire (la première étant décroissante, les produits \deleted{\unsure{qui interviennent par}} des $\frac1{\rho_i}$ ég. au produit des $\frac1{\sigma_i}$, et enfin $\frac1{\rho_n}\ldots\frac1{\rho_i}\leq\frac1{\sigma_n}\ldots\frac1{\sigma_i}$ pour $1\leq M\leq n$), d'où l'inégalité (10).

% archive sheet number: 86
